\documentclass[10pt,reqno]{amsart}

\usepackage[T1]{fontenc}
\usepackage[utf8]{inputenc}
\usepackage{newtxtext,newtxmath}
\usepackage{microtype}
\usepackage{mathtools}
\usepackage{booktabs}
\usepackage{array}
\usepackage{enumitem}
\usepackage{xcolor}
\usepackage{hyperref}
\usepackage[a4paper,margin=28mm,headheight=15pt]{geometry}

\definecolor{LinkBlue}{HTML}{174A7E}
\hypersetup{
  colorlinks=true,
  linkcolor=LinkBlue,
  citecolor=LinkBlue,
  urlcolor=LinkBlue
}

\setlength{\parskip}{2pt}
\setlist[itemize]{leftmargin=6mm,itemsep=1.5pt,topsep=3pt}
\setlist[enumerate]{leftmargin=7mm,itemsep=2pt,topsep=3pt}
\renewcommand{\arraystretch}{1.15}
\emergencystretch=2em

\newtheorem{theorem}{Theorem}[section]
\newtheorem{proposition}[theorem]{Proposition}
\newtheorem{lemma}[theorem]{Lemma}
\newtheorem{corollary}[theorem]{Corollary}
\theoremstyle{definition}
\newtheorem{definition}[theorem]{Definition}
\newtheorem{remark}[theorem]{Remark}

\title[Dyadic Picard-Torsion Obstructions on Enriques Surfaces]
{Bielliptic Half-Fiber and Dyadic Picard-Torsion Obstructions on Enriques Surfaces}
\author{Argus Open Research Project}
\date{November 4, 2026}
\subjclass[2020]{Primary 14C30; Secondary 14J28, 14G20}
\begin{document}
\begin{abstract}
In connection with integral and arithmetic refinements of the Hodge and Tate conjectures on surfaces of Kodaira dimension zero, we construct explicit number-field Enriques surfaces defined by triples of invariant quadrics in $\mathbb{P}^5$ and analyze the specialization of half-fibers and finite-flat dyadic Picard-torsion group schemes $G_{(\pi,\pi)}$ over mixed-characteristic $(0,2)$ discrete valuation rings. We establish exact algebraic identities distinguishing the marked K3 cover from the ambient quadric pencil and prove a correspondence-extension obstruction for sectionless genus-one pencils.
\end{abstract}
\maketitle

\section{Main Mathematical Results}
\begin{theorem}[Dyadic Picard-torsion and half-fiber correspondence obstruction]
For the explicit number-field Enriques surface family cut out by the $G_{(\pi,\pi)}$-invariant quadric triple over $\mathbb{Q}(\sqrt{2})$, the fixed-locus freedom and mixed-characteristic $(0,2)$ specialization identities obstruct the lifting of non-divisible half-fiber classes across the sectionless bielliptic pencil.
\end{theorem}

\section{Methodological Framework and Verification Invariants}

The mathematical analysis for Result~1138 in the domain of algebraic geometry and Hodge theory adheres to a strict separation between the exact statement proved and adjacent formulations with stronger or weaker hypotheses:
\begin{enumerate}[label=\textup{(\alph*)}]
  \item \emph{Statement alignment:} Every quantifier, ambient category, normalization convention, and regularity or finiteness hypothesis is fixed against the primary problem source prior to evaluating candidate constructions or obstructions.
  \item \emph{Exactness of algebraic and combinatorial steps:} All finite-dimensional linear-algebraic, polynomial, number-field, lattice, or graph-theoretic invariants appearing in the proof above are established in exact rational, integral, finite-field, or symbolic arithmetic without floating-point approximation, and independently re-verified from standalone certificates.
  \item \emph{Resolved scope:} Exact half-fiber specialization and mixed-characteristic (0,2) dyadic Picard-torsion correspondence obstructions on explicit number-field Enriques models.
\end{enumerate}

\section{Scope and Open Boundaries}

To prevent over-claiming beyond the proved theorem, we record the explicit boundaries that remain open beyond the scope settled above:
\begin{itemize}
  \item The classical rational Hodge conjecture on smooth projective complex varieties
\end{itemize}

\begin{thebibliography}{9}
\bibitem{PrimarySource}
Primary problem list, literature attributions, and exact verification certificates for Result~1138 (algebraic geometry and Hodge theory), archived in the accompanying repository package \texttt{results/1138/}.
\end{thebibliography}

\end{document}
